% sec_01_front.tex -- fragment included by memo.tex. Title, author, date and the
% status block live in the master file.

\begin{abstract}
\noindent
This memo asks under what product, demand, regulatory and operating conditions regional
manufacturing of sterile injectables could reduce shortages more effectively than safety
stock, dual sourcing, or additional centralized capacity. It compares those options on two
non-controlled small-molecule presentations: sodium bicarbonate 8.4 percent in a 50 mL vial,
in prolonged shortage, and norepinephrine bitartrate 1 mg/mL in a 4 mL vial, resolved. The
method is a protocol frozen before any result, a deterministic screen reconciled cell for
cell against the source spreadsheet, a daily discrete-event simulation in which every
strategy faces the same simulated world, optimization of each strategy to one pre-specified
service target, and an ablation harness over named failure mechanisms. The answer is
conditional and largely negative: what decides the outcome is capacity relative to demand and
correlated failure across nominally separate sites, and neither is improved by distributing
manufacture as such. The study retracted three of its own headline conclusions after
correcting defects in its own model, and that history is reported here rather than removed.
\end{abstract}

\section{The problem, the two products, and the question}
\label{sec:question}

\subsection{What the record supports}
\label{sec:record}

United States sterile injectable shortages are counted by two registries that do not count the
same object. The ASHP and University of Utah series recorded 323 active shortages in the first
quarter of 2024, the highest since tracking began in 2001. The FDA CDER database, which is the
statutory list, carried 102 as of 31 July 2024 in the count the Government Accountability
Office records. Both are correct on their own definitions and the gap is roughly threefold.
This package flagged that its own earlier market research quoted the 323 figure without saying
it was the ASHP count, which overstates the FDA-defined problem. The distinction is
operational, because the FDA list is the legal object that gates the 503B exception, the CMS
buffer-stock payment and most contracting language, so a claim to reduce shortage days is
measured against that list.

Three features of the record bound what a manufacturing intervention can do. Demand, not plant
failure, is the dominant reported driver: increased demand contributed to 66 of the 102
shortages on the FDA list at that date. That driver is also the least observable, because
manufacturers must report an interruption caused by increased demand but need not report an
increase in demand that does not interrupt production, and FDA has asked Congress for that
authority. And concentration risk is not only offshore: the largest recent supply event
followed a hurricane closing a single domestic plant in Marion, North Carolina, described by
USP as having made roughly 60 percent of US intravenous solution supply.

The prior-art review behind this study covered fourteen solution families over 798 source rows,
651 fetched and quoted, with citation audits on 2 and 3 September 2026. The audits found no
fabricated source, company, patent, contract or DOI, and found instead a recurring defect of
misattribution, load-bearing numbers cited to sources that do not contain them, with thirty-six
corrections in four families alone. The structural finding is that the space is occupied. A
centralized plant built specifically for shortage injectables is already implemented rather
than proposed, and shortage detection, risk scoring, ingredient mapping, essential-medicines
curation and committed-volume contracting each have a funded incumbent with better data than a
new entrant would have. What is unsettled is the economics and operations of a distributed
alternative, which is why this is a feasibility question rather than a novelty question, and
why it is built so the answer may be negative.

\subsection{Two products, from a screen that selects nothing}
\label{sec:products}

The protocol requires a dated longlist of at least five presentations from frozen FDA and ASHP
snapshots including one resolved or recurrent comparator, then hard gates, then weighted
scoring of survivors. Six were screened on 2 September 2026 against snapshots whose hashes are
recorded in \texttt{data/product\_dossiers/screen\_manifest.json}.

\begin{table}[htbp]
\centering
\caption{Product screen, 2 September 2026, against frozen official snapshots. Shortage rows
from the FDA drug shortage export (1,635 rows). No candidate passed every gate, because gates
PG5 to PG8 cannot be evaluated from public data.}
\label{tab:screen}
\footnotesize
\setlength{\tabcolsep}{4pt}
\begin{tabular}{@{}llllc@{}}
\toprule
Presentation & Protocol role & Binding gate result & FDA rows (current) & Outcome \\
\midrule
Sodium bicarbonate 8.4\%, 50 mL vial & case candidate & PG1 to PG3 PASS & 19 (15) & UNCERTAIN \\
Norepinephrine bitartrate 1 mg/mL, 4 mL & resolved comparator & PG1 to PG3 PASS & 0 (0) & UNCERTAIN \\
Furosemide 10 mg/mL vial & reserve & PG1 UNCERTAIN & 33 (30) & UNCERTAIN \\
Sterile water for injection, vial & scale question & PG1 UNCERTAIN, PG8 mismatch & 21 (21) & UNCERTAIN \\
Acyclovir sodium 50 mg/mL vial & reserve & PG1 UNCERTAIN & 0 (0) & UNCERTAIN \\
Acetazolamide sodium 500 mg vial & out-of-archetype & PG2 FAIL, lyophilized & 0 (0) & FAIL \\
\bottomrule
\end{tabular}
\end{table}

The two carried forward are sodium bicarbonate injection 8.4 percent in a 50 mL single-dose
vial and norepinephrine bitartrate 1 mg/mL in a 4 mL single-dose vial, chosen for the contrast
the protocol asks for. Sodium bicarbonate is the prolonged current shortage: 19 rows in the FDA
export, 11 matching the 8.4 percent presentation, 15 current, earliest posting 1 March 2017,
stated reasons demand increase, discontinuation and other. Norepinephrine is the resolved
comparator: zero rows in the FDA export and the openFDA mirror on the same date, and one ASHP
bulletin ever, opened 9 February 2017 and resolved 7 June 2022. Both are non-controlled
small-molecule aqueous solutions in standard vials, the archetype the modelled process train
represents, and both clear the identity, manufacturing-fit and control-status gates that
acetazolamide fails on lyophilization. Their supplier structures differ, 13 approved holders
against 23, and under the illustrative inputs one network is capacity-short by construction and
the other is not. That contrast is the point of the pair: a conclusion holding on only one of
them is a conclusion about a product rather than about an architecture.

Four qualifications attach. The screen selects nothing, because gates PG5 to PG8, covering
demand observability, clinical substitutability, input feasibility and economic relevance,
cannot be evaluated from public data, so every survivor ends UNCERTAIN and the protocol's rule
is that UNCERTAIN blocks final selection. The protocol's step 5 requires a hospital pharmacy
leader and a sterile-manufacturing expert to challenge the selection independently, and that
has not happened. The selection record contains zero finalized presentations and both product
configurations carry a header saying they are provisional. And the source workbook labelled
norepinephrine ``resolved/recurrent'' where the snapshot does not support the recurrent half.
These are working presentations for a model, used as such.

\subsection{The question}
\label{sec:central-question}

\begin{quote}
Under what product, demand, regulatory and operating conditions could regional manufacturing
reduce sterile-injectable shortages more effectively than safety stock, dual sourcing, or
additional centralized capacity?
\end{quote}

The words that carry weight are ``under what conditions''. The question is not whether
distributed manufacturing is a good idea but which region of the product, demand, regulatory
and operating space, if any, contains a distributed design that beats the alternatives on cost
at a matched service level. A negative answer, or one holding only inside a narrow and clearly
stated region, is a legitimate result and the study was built to permit one. The comparison set
is fixed in the protocol as the eight strategies S0 to S7 of Table~\ref{tab:architectures},
which implement the four comparator classes the specification names and add centralized
expansion and a 503B responder; twelve further architectures were configured later as S8 to S19
and are compared on the same terms. Quality and regulatory requirements enter as hard
constraints throughout. They are never relaxed to improve an economic result, and an ineligible
strategy is excluded rather than charged a penalty an optimizer could pay through.

\section{Method}
\label{sec:method}

The protocol was written and frozen first, as \texttt{protocol/protocol.yaml} version 1.0.0,
converted on 1 September 2026 from the author's protocol document and the source workbook. The
service target was fixed before any simulation result existed: mean fill rate at least 0.99,
and probability at least 0.90 that a run's own fill rate reaches 0.99, with a shortage day
defined as a day whose fulfilled fraction falls below 0.95. Those values are recorded in
revision R001 on the day of the freeze, with the note that no simulation result existed at that
date. Eleven revisions have been logged since, ten flagged decision-relevant, each carrying a
date, section, change, rationale, author, approval state and affected files
(Table~\ref{tab:revisions}). Most are corrections the study made to its own model after finding
a defect in it. None has been signed off, so the definition-of-finished gate for protocol freeze
fails and is recorded as failing.

\begin{table}[htbp]
\centering
\caption{The eleven logged protocol revisions (\texttt{protocol/revisions.csv}). ``DR'' marks
one flagged decision-relevant. None has been signed off.}
\label{tab:revisions}
\footnotesize
\setlength{\tabcolsep}{5pt}
\begin{tabular}{@{}llll@{}}
\toprule
Id & Date & Subject & DR \\
\midrule
R001 & 2026-09-01 & Service thresholds pre-specified before any result & yes \\
R002 & 2026-09-01 & Regulatory gate set adopted, every status UNCERTAIN & yes \\
R003 & 2026-09-02 & Raw-material stock carries the inventory carrying rate & yes \\
R004 & 2026-09-03 & Expansion lead time, cost accrual, opening cohorts, reserved lines & yes \\
R005 & 2026-09-03 & Reorder point, node placement, utilization denominator, error bars & no \\
R006 & 2026-09-03 & Twelve further architectures configured as S8 to S19, gate-mapped & yes \\
R007 & 2026-09-03 & Supplier capacity as a fraction, reserved capacity screened, gate G16 & yes \\
R008 & 2026-09-03 & Commissioning clock starts at the first measured day & yes \\
R009 & 2026-09-05 & Opening inventory is purchased and charged, not inherited free & yes \\
R010 & 2026-09-05 & Every comparator searches one common inventory space & yes \\
R011 & 2026-09-05 & Pipeline cover, region weights, readiness clock, contract priority & yes \\
\bottomrule
\end{tabular}
\end{table}

The study inherited an Excel model and reproduces it rather than replacing it silently.
\texttt{deterministic.run\_screen} in replica mode recomputes the screen for all sixteen rows
from the same base inputs, and \texttt{reconcile\_with\_workbook} reports any absolute
difference above $10^{-6}$ for currency or $10^{-9}$ for fractions. All 448 cached values
reconcile, the workbook snapshot is pinned by SHA-256 hash, and the reconciliation is a
regression test rather than a one-time exercise. A separate corrected mode applies the defects
the workbook audit found, including a delivered-unit denominator and a capacity cap, and
reports its deltas separately, so a correction cannot be misread as a reconciliation failure.
The screen is a verification device and produces no service verdict, because a deterministic
screen cannot produce one.

Service and cost come from a daily discrete-event simulation of sites, suppliers, regions and
lanes. Each run generates one exogenous world for a universal roster of entities, and every
strategy in that run consumes the subset it instantiates, so paired strategies see identical
demand, shocks, site failures, supplier disruptions, common-cause events, shortage-list
transitions and transit noise. This is the study's first fair-comparison rule and it is enforced
by test. Random draws are keyed by run index, named stream, a BLAKE2b hash of the entity
identifier and an ordinal, so adding an entity to a design never perturbs another entity's
draws, which is what makes a comparison between two designs a comparison of the designs.

No strategy is compared at its default settings. For each product and strategy the optimizer
reads regulatory eligibility from the gates first, so an excluded strategy is never optimized
and a strategy whose gates are unresolved is optimized for information and labelled as such. It
searches the strategy's declared design space on a common random-number batch, where feasible
means the mean fill rate reaches the target and the share of runs reaching it clears the tail
confidence, refines the best feasible point by coordinate search, and re-evaluates the incumbent
at full sample size with its Monte Carlo standard error. If no point is feasible the record says
so and reports the closest fill rate. Nothing pays a penalty to become feasible.

Attribution comes from an ablation harness of 22 named failure mechanisms, each a set of
parameter overrides applied through the same provenance-carrying records the rest of the study
uses, so a switched-off mechanism is a stated value visible in the run manifest. For each
product and strategy the harness runs a base configuration, a leave-one-out configuration per
mechanism, an all-off configuration, and an add-one-in configuration per mechanism.
Leave-one-out and add-one-in are the two ends of the Shapley averaging order, so the harness
reports the bracket between them rather than a single attributed share, and reports the Monte
Carlo standard error of every mean, on the rule that a difference smaller than the standard
errors it is built from is not evidence of anything.

The harness turned out to be a defect detector as much as an attribution tool, because a
mechanism whose measured sign is wrong is usually a bug. That is how most of the twenty-four
model defects in \texttt{bottleneck\_decomposition.md} section 7 were found, and it is why this
memo carries a correction history. After defects NEW-1, MD-3, MD-11 and MD-12 were fixed under
revisions R009 to R011 and every incumbent design was re-evaluated at full resolution without
re-searching, three statements the study had presented as headline conclusions were falsified by
its own corrected model and withdrawn; Section~\ref{sec:limitations} reports them. Three limits
bound everything that follows, beyond the standing qualification on page 1. The ablation
measures only mechanisms the engine represents, so the absence of a mechanism is not evidence
that it does not matter; known omissions include multi-product pooling, campaign scheduling
across products, and a second inventory stage for postponement. The engine models one product
per network over four equal regions. And the entire input set is illustrative, so the ranking of
mechanisms, like every cost and service number here, is a property of those values. Where a
result would change under a different input, the memo reports the condition rather than the
number.
